% Generated draft: requires mathematical and source review.
\documentclass[11pt]{article}
\usepackage[margin=25mm]{geometry}
\usepackage{fontspec}
\setmainfont{Latin Modern Roman}
\usepackage{amsmath,amssymb,mathtools,unicode-math}
\setmathfont{Latin Modern Math}
\usepackage{enumitem,needspace,xurl,hyperref}
\setlength{\parindent}{0pt}
\setlength{\parskip}{5pt}
\setlength{\emergencystretch}{3em}
\newcommand{\N}{\mathbb{N}}
\newcommand{\Z}{\mathbb{Z}}
\newcommand{\Q}{\mathbb{Q}}
\newcommand{\R}{\mathbb{R}}
\newcommand{\C}{\mathbb{C}}
\newcommand{\F}{\mathbb{F}}
\newcommand{\A}{\mathbb{A}}
\newcommand{\PP}{\mathbb{P}}
\newcommand{\E}{\mathbb{E}}
\newcommand{\eps}{\varepsilon}
\newcommand{\dd}{\,\mathrm{d}}
\begin{document}
\section*{UnsolvedMath ID: 30001358}
\textbf{Source problem code:} catalogue entry
\subsection{Definitions and mathematical statement}
Equip a smooth connected embedded surface \(M\subset\R^3\) with the induced Riemannian metric \(g=\iota^*g_{\mathrm{Euc}}\), where \(\iota:M\hookrightarrow\R^3\) is inclusion. Completeness means completeness of the induced Riemannian distance, and minimality means that the mean curvature vector vanishes everywhere. An intrinsic isometry is a diffeomorphism \(\phi:M\to M\) satisfying \(\phi^*g=g\). A Euclidean isometry has the form \(A(x)=Qx+b\), where \(Q\in O(3)\) and \(b\in\R^3\).

The extension clause of the supplied question asks: for every connected complete embedded minimal surface \(M\subset\R^3\) and every intrinsic isometry \(\phi:M\to M\), do there exist \(Q\in O(3)\) and \(b\in\R^3\) such that
\[
\iota(\phi(p))=Q\iota(p)+b\qquad\text{for every }p\in M?
\]
This formulation uses the usual convention that a surface has no boundary. The record does not explicitly specify the boundary convention; whether surfaces with boundary are intended requires clarification. No properness, orientability, finite-genus, or periodicity hypothesis is supplied.

The source additionally asks for ``minimal rigidity unless \(M\) is a helicoid.'' Here a helicoid means a surface congruent to
\[
H_a=\{(s\cos t,s\sin t,at):s,t\in\R\},\qquad a\in\R\setminus\{0\}.
\]
The source does not define minimal rigidity or its comparison class. If it means rigidity among all isometric minimal immersions of the same Riemannian surface, the additional assertion would be: whenever \(M\) is not a helicoid, every smooth minimal immersion \(F:M\to\R^3\) satisfying \(F^*g_{\mathrm{Euc}}=g\) satisfies
\[
F=A\circ\iota
\]
for some Euclidean isometry \(A\). This is an explicitly conditional interpretation requiring source verification: the supplied record does not determine whether the comparison maps must be immersions, embeddings, or members of a deformation family, or whether reparametrization is allowed in the congruence relation. Consequently, the full rigidity clause cannot yet be assigned a unique precise formulation. The helicoid exception is attached to this additional rigidity clause, not to the stated intrinsic-isometry extension clause.
\subsection{Short English statement}
The question asks whether every intrinsic symmetry of a connected complete embedded minimal surface in \(\R^3\) comes from a Euclidean motion. It also asks for minimal rigidity away from helicoids, but the intended meaning of rigidity and the boundary convention require clarification. The supplied status is partially solved; this draft does not independently verify that assessment.
\subsection{Sources}
\begin{itemize}
\item Supplied local UnsolvedMath registry record, permanent ID \(30001358\), titled ``Rigidity of Complete Embedded Minimal Surfaces,'' in the Oberwolfach Reports Open Problems collection; exact source problem code \texttt{catalogue entry}. Project registry locations are \texttt{data/unsolvedmath/problems.json} and the supplementary \texttt{catalog/problems.json}; no independent registry inspection is claimed.
\item Supplied source citation: \emph{Arbeitsgemeinschaft: Minimal Surfaces (2010)}, Oberwolfach Reports, DOI 10.4171/owr/2009/45, pp. 2545--2584. \url{https://doi.org/10.4171/owr/2009/45}. The record gives proposed year 2009 and names no proposer. Its background describes Oberwolfach Reports as documenting mathematical workshops and open-problem discussions.
\item Supplied literature reference: Meeks--P\'erez, minimal-surface survey. \url{https://www.intlpress.com/site/pub/files/_fulltext/journals/sdg/2016/0021/0001/SDG-2016-0021-0001-a006.pdf}.
\item Supplied literature reference: Collin--Kusner--Meeks--Rosenberg, end restrictions. \url{https://arxiv.org/abs/math/0401267}.
\item Supplied literature reference: original OWR report. \url{https://doi.org/10.4171/owr/2009/45}.
\item Literature assessment dated exactly \texttt{2026-08-21}: the supplied curation labels the problem partially solved and reports that the Isometry Conjecture is proved in several substantial classes, including finite-genus and periodic cases, while the general infinite-genus case remains open. This is supplied dated curation derived from the OWR record, not independently verified here; its time-sensitive status and cited claims require independent verification.
\item The supplied statement-status label is \texttt{corrected\_verified}. Its accompanying verification note says only that TeX notation was enclosed in mathematical delimiters without changing mathematical content; it does not establish mathematical correctness or source verification.
\end{itemize}
\end{document}
